\textbf{The property.} An update is accepted when its operand's cell is \emph{unique at the site}: every other holder of that cell is dead at that point. A holder is dead when no later read goes through it. A read that comes \emph{before} the update is therefore never sharing, however many there are; a read that comes \emph{after}, through any path the checker can see, is. This is Boyland's alias burying (2001 §3.2, p. 13: aliases may exist “simply … all be dead”) and Rust's non-lexical borrow (RFC 2094, “Liveness”), applied to a language with no references, and it is exactly what a plain array needs: the write is invisible iff nobody reads the old version afterwards.
